% Two-pass linear-time tree solver: (a) bottom-up subtree injection sum for edge currents, (b) top-down potential recovery.
\begin{tikzpicture}[x=1mm,y=1mm,font=\figurefont\fontsize{8}{10}\selectfont,>=Stealth,
  vertex/.style={circle,draw=FigureBlue,fill=FigureBlueFill,inner sep=0pt,
    minimum size=5mm,line width=.55pt},
  rootv/.style={circle,draw=FigureGreen,fill=FigureGreenFill,inner sep=0pt,
    minimum size=5.2mm,line width=.68pt},
  upflow/.style={->,draw=FigureBlue,line width=.8pt},
  downflow/.style={->,draw=FigureGreen,line width=.8pt},
  note/.style={font=\FigureNoteFont\fontsize{7}{8.5}\selectfont,text=FigureInk,inner sep=1pt},
  panel/.style={font=\figurefont\fontsize{7.5}{9}\selectfont,text=FigureInk}]
  \definecolor{FigureBlue}{HTML}{7188A3}
  \definecolor{FigureBlueFill}{HTML}{E5EAF0}
  \definecolor{FigureGreen}{HTML}{829C8B}
  \definecolor{FigureGreenFill}{HTML}{E8EEE9}
  \definecolor{FigureYellow}{HTML}{B4A16F}
  \definecolor{FigureYellowFill}{HTML}{F1ECDD}
  \definecolor{FigureRed}{HTML}{A7787D}
  \definecolor{FigureInk}{HTML}{3B4252}
  \definecolor{FigureMuted}{HTML}{6B7280}
  \definecolor{FigureStroke}{HTML}{4C4D4F}

  % Panel (a): Bottom-up pass (leaves to root 4)
  \begin{scope}[xshift=0mm]
    \node[rootv]  (r4) at (18,22) {$4$};
    \node[vertex] (v3) at (9,11)  {$3$};
    \node[vertex] (v5) at (27,11) {$5$};
    \node[vertex] (v1) at (2,0)   {$1$};
    \node[vertex] (v2) at (16,0)  {$2$};

    \node[note,above=1mm] at (r4.north) {$b_4=-2$};
    \node[note,left=1mm]  at (v3.west)  {$b_3=-1$};
    \node[note,right=1mm] at (v5.east)  {$b_5=+1$};
    \node[note,below=1mm] at (v1.south) {$b_1=+1$};
    \node[note,below=1mm] at (v2.south) {$b_2=+1$};

    \draw[upflow] (v1) -- node[left,note,text=FigureBlue] {$s_1=1$} (v3);
    \draw[upflow] (v2) -- node[right,note,text=FigureBlue] {$s_2=1$} (v3);
    \draw[upflow] (v3) -- node[above left,note,text=FigureBlue] {$s_3=1$} (r4);
    \draw[upflow] (v5) -- node[above right,note,text=FigureBlue] {$s_5=1$} (r4);

    \node[panel] at (15,-10) {(a) 第一遍：自底向上汇总得边电流 $s_i$};
  \end{scope}

  % Panel (b): Top-down pass (root 4 to leaves, unit resistance r_e = 1)
  \begin{scope}[xshift=58mm]
    \node[rootv]  (br4) at (18,22) {$4$};
    \node[vertex] (bv3) at (9,11)  {$3$};
    \node[vertex] (bv5) at (27,11) {$5$};
    \node[vertex] (bv1) at (2,0)   {$1$};
    \node[vertex] (bv2) at (16,0)  {$2$};

    \node[note,above=1mm,text=FigureGreen] at (br4.north) {$v_4=0$（根）};
    \node[note,left=1mm,text=FigureInk]    at (bv3.west)  {$v_3=1$};
    \node[note,right=1mm,text=FigureInk]   at (bv5.east)  {$v_5=1$};
    \node[note,below=1mm,text=FigureInk]   at (bv1.south) {$v_1=2$};
    \node[note,below=1mm,text=FigureInk]   at (bv2.south) {$v_2=2$};

    \draw[downflow] (br4) -- node[above left,note,text=FigureGreen] {$+r_{34}s_3$} (bv3);
    \draw[downflow] (br4) -- node[above right,note,text=FigureGreen] {$+r_{54}s_5$} (bv5);
    \draw[downflow] (bv3) -- node[left,note,text=FigureGreen] {$+r_{13}s_1$} (bv1);
    \draw[downflow] (bv3) -- node[right,note,text=FigureGreen] {$+r_{23}s_2$} (bv2);

    \node[panel] at (15,-10) {(b) 第二遍：自顶向下递推得电势 $v_i$};
  \end{scope}
\end{tikzpicture}
